
%%%%%%%%%%%%%%%%%%%%%%%%%%%%%%%%%%%%%%%%%%%%%%%%%%%%%%%%%%%%%%%%%%%%%%%%%%%%
%%%%%%%%%%%%%%% ESTO ES PARA LA TABLA DE CONTENIDOS (INDICE)%%%%%%%%%%%%%%%%


\setcounter{secnumdepth}{3} %niveles en capitulos hasta (1.1.1)
\setcounter{tocdepth}{3} %niveles en indice hasta (1.1)


\addtocontents{toc}{~\hfill\small{Página}\par} %añade "pagina" a tabla de contanidos
\addtocontents{toc}{\vspace{2pt} \hrule \vspace{5mm} \par}

%\renewcommand\contentsname{Índice de contenidos}
\tableofcontents
\listoffigures % indice de figuras
\listoftables % indice de tablas
\lstlistoflistings % indice de listados: comentarlo si la memoria no incluye
% bloques de código (listings); si no, genera una página vacía titulada
% "LISTADOS"



%---------------------------------------------------------------------------
% comienzo de relacion abreviaturas y/o acrónimos
%---------------------------------------------------------------------------


%\addcontentsline{toc}{section}{ABREVIATURAS}
\clearpage
\vspace{0.2cm}
\section*{ABREVIATURAS}
% Tabla de ejemplo. Sustitúyela por los acrónimos y abreviaturas que
% realmente uses en tu memoria, ordenados alfabéticamente. Cada acrónimo
% debe definirse la primera vez que se usa en el texto, ver capítulo
% "Ejemplos prácticos de la plantilla".
% La tabla ocupa el ancho del texto con columnas elásticas (Y, definidas en
% main.tex) en lugar de \resizebox, que encogía la letra: las descripciones
% largas hacen salto de línea. Se lee por columnas: primero la pareja de la
% izquierda y después la de la derecha.
\noindent
\begin{tabularx}{\textwidth}{|Y{0.4600}|Y{1.5400}|Y{0.4600}|Y{1.5400}|}
\hline
API & \textit{Application Programming Interface} & TFG & Trabajo Fin de Grado \\
ESI & Escuela Superior de Ingeniería & TFM & Trabajo Fin de Máster \\
IDE & \textit{Integrated Development Environment} & UAL & Universidad de Almería \\
IEEE & \textit{Institute of Electrical and Electronics Engineers} & UML & \textit{Unified Modeling Language} \\
PDF & \textit{Portable Document Format} & URL & \textit{Uniform Resource Locator} \\
\hline
\end{tabularx}
%\addcontentsline{toc}{section}{ABREVIATURAS}

